% sec4_5.tex — Section 4.5: Special Cases of Multiple-Equation GMM: FIVE, 3SLS, and SUR

\section{Special Cases of Multiple-Equation GMM: FIVE, 3SLS, and SUR}
\label{sec:4.5}

\subsection*{Conditional Homoskedasticity}

Under conditional homoskedasticity, we can easily derive a number of prominent
estimators as special cases of multiple-equation GMM. The multiple-equation version
of conditional homoskedasticity is

\begin{assumption}[conditional homoskedasticity]
\label{ass:4.7}
\[
\E(\varepsilon_{im}\, \varepsilon_{ih} \mid \mathbf{x}_{im}, \mathbf{x}_{ih}) = \sigma_{mh}
\]
\textit{for all $m, h = 1, 2, \ldots, M$.}
\end{assumption}

Then the unconditional cross moment $\E(\varepsilon_{im}\, \varepsilon_{ih})$ equals
$\sigma_{mh}$ by the Law of Total Expectations. You should have no trouble showing that
\begin{align}
(m, h) \text{ block of } \E(\mathbf{g}_i \mathbf{g}_i') \text{ given in \eqref{eq:4.1.11}}
&= \E(\varepsilon_{im}\, \varepsilon_{ih}\, \mathbf{x}_{im} \mathbf{x}_{ih}') \notag \\
&= \E[\E(\varepsilon_{im}\, \varepsilon_{ih}\, \mathbf{x}_{im} \mathbf{x}_{ih}'
\mid \mathbf{x}_{im}, \mathbf{x}_{ih})]
&& \text{(by the Law of Total Expectations)} \notag \\
&= \E[\E(\varepsilon_{im}\, \varepsilon_{ih} \mid \mathbf{x}_{im}, \mathbf{x}_{ih})
\mathbf{x}_{im} \mathbf{x}_{ih}']
&& \text{(by linearity of conditional expectations)} \notag \\
&= \E[\sigma_{mh}\, \mathbf{x}_{im} \mathbf{x}_{ih}']
&& \text{(by conditional homoskedasticity)} \notag \\
&= \sigma_{mh}\, \E(\mathbf{x}_{im} \mathbf{x}_{ih}')
&& \text{(by linearity of expectations).}
\label{eq:4.5.1}
\end{align}

Thus, the $\mathbf{S}$ in~\eqref{eq:4.1.11} can be written as
\begin{equation}
\mathbf{S} = \begin{bmatrix}
\sigma_{11}\, \E(\mathbf{x}_{i1} \mathbf{x}_{i1}') & \cdots &
\sigma_{1M}\, \E(\mathbf{x}_{i1} \mathbf{x}_{iM}') \\
\vdots & & \vdots \\
\sigma_{M1}\, \E(\mathbf{x}_{iM} \mathbf{x}_{i1}') & \cdots &
\sigma_{MM}\, \E(\mathbf{x}_{iM} \mathbf{x}_{iM}')
\end{bmatrix}.
\label{eq:4.5.2}
\end{equation}

Since by Assumption~4.5 $\mathbf{S}$ is finite, this decomposition implies that
$\E(\mathbf{x}_{im} \mathbf{x}_{ih}')$ exists and is finite for all
$m, h$ $(= 1, 2, \ldots, M)$.

\subsection*{Full-Information Instrumental Variables Efficient (FIVE)}

An estimator of $\mathbf{S}$ exploiting the structure of fourth moments shown
in~\eqref{eq:4.5.2} is
\begin{equation}
\widehat{\mathbf{S}} = \begin{bmatrix}
\hat{\sigma}_{11} \cdot \Bigl(\dfrac{1}{n} \displaystyle\sum_{i=1}^{n}
\mathbf{x}_{i1} \mathbf{x}_{i1}'\Bigr) & \cdots &
\hat{\sigma}_{1M} \cdot \Bigl(\dfrac{1}{n} \displaystyle\sum_{i=1}^{n}
\mathbf{x}_{i1} \mathbf{x}_{iM}'\Bigr) \\[8pt]
\vdots & & \vdots \\[4pt]
\hat{\sigma}_{M1} \cdot \Bigl(\dfrac{1}{n} \displaystyle\sum_{i=1}^{n}
\mathbf{x}_{iM} \mathbf{x}_{i1}'\Bigr) & \cdots &
\hat{\sigma}_{MM} \cdot \Bigl(\dfrac{1}{n} \displaystyle\sum_{i=1}^{n}
\mathbf{x}_{iM} \mathbf{x}_{iM}'\Bigr)
\end{bmatrix},
\label{eq:4.5.3}
\end{equation}
where, for some consistent estimator $\hat{\boldsymbol{\delta}}_m$ of $\boldsymbol{\delta}$,
\begin{equation}
\hat{\sigma}_{mh} \equiv \frac{1}{n} \sum_{i=1}^{n} \hat{\varepsilon}_{im}\, \hat{\varepsilon}_{ih},
\quad
\hat{\varepsilon}_{im} \equiv y_{im} - \mathbf{z}_{im}' \hat{\boldsymbol{\delta}}_m
\quad (m, h = 1, 2, \ldots, M).
\label{eq:4.5.4}
\end{equation}

By Proposition~4.1, $\hat{\sigma}_{mh} \pto \sigma_{mh}$ provided (in addition to
Assumptions~4.1 and~4.2) that $\E(\mathbf{z}_{im} \mathbf{z}_{ih}')$ is finite.
By ergodic stationarity,
$\frac{1}{n} \sum_i \mathbf{x}_{im} \mathbf{x}_{ih}'$ converges in probability to
$\E(\mathbf{x}_{im} \mathbf{x}_{ih}')$, which, as just noted, exists and is finite.
Therefore,~\eqref{eq:4.5.3} is consistent for $\mathbf{S}$, without the fourth-moment
assumption (Assumption~4.6).

The \textbf{FIVE estimator} of $\boldsymbol{\delta}$, denoted
$\hat{\boldsymbol{\delta}}_{\mathrm{FIVE}}$, is\footnote{The estimator is due to
Brundy and Jorgenson (1971).}
\[
\hat{\boldsymbol{\delta}}_{\mathrm{FIVE}}
\equiv \hat{\boldsymbol{\delta}}(\widehat{\mathbf{S}}^{-1})
\]
with $\widehat{\mathbf{S}}$ given by~\eqref{eq:4.5.3} to exploit conditional
homoskedasticity.\footnote{This assumption is needed for $\hat{\sigma}_{mh}$ to be
consistent. See Proposition~4.1.}
Being a multiple-equation GMM for some choice of $\hat{\mathbf{W}}$, it is consistent
and asymptotically normal by (the multiple-equation adaptation of)
Proposition~3.1. Because $\widehat{\mathbf{S}}$ is consistent for $\mathbf{S}$,
the estimator is efficient by Proposition~3.5. Thus, we have proved

\begin{proposition}[large-sample properties of FIVE]
\label{prop:4.4}
\textit{Suppose Assumptions~4.1--4.5 and~4.7 hold. Suppose, furthermore, that\/
$\E(\mathbf{z}_{im} \mathbf{z}_{ih}')$ exists and is finite for all
$m, h$ $(= 1, 2, \ldots, M)$.
Let\/ $\mathbf{S}$ and\/ $\widehat{\mathbf{S}}$ be as in~\eqref{eq:4.5.2}
and~\eqref{eq:4.5.3}, respectively. Then}
\begin{enumerate}[label=(\alph*)]
\item $\widehat{\mathbf{S}} \pto \mathbf{S}$;
\item $\hat{\boldsymbol{\delta}}_{\mathrm{FIVE}}$, \textit{defined as\/
$\hat{\boldsymbol{\delta}}(\widehat{\mathbf{S}}^{-1})$, is consistent,
asymptotically normal, and efficient, with\/
$\avar(\hat{\boldsymbol{\delta}}_{\mathrm{FIVE}})$ given by~\eqref{eq:4.3.3};}
\item \textit{The estimated asymptotic variance given in~\eqref{eq:4.3.4} is
consistent for\/ $\avar(\hat{\boldsymbol{\delta}}_{\mathrm{FIVE}})$;}
\item \textit{(Sargan's statistic)}
\[
J(\hat{\boldsymbol{\delta}}_{\mathrm{FIVE}},\, \widehat{\mathbf{S}}^{-1})
\equiv n \cdot \mathbf{g}_n(\hat{\boldsymbol{\delta}}_{\mathrm{FIVE}})'
\widehat{\mathbf{S}}^{-1}\, \mathbf{g}_n(\hat{\boldsymbol{\delta}}_{\mathrm{FIVE}})
\dto \chi^2\!\Bigl(\sum_m (K_m - L_m)\Bigr),
\]
\textit{where\/ $\mathbf{g}_n(\cdot)$ is given in~\eqref{eq:4.2.1}.}
\end{enumerate}
\end{proposition}

Usually, the initial estimator $\hat{\boldsymbol{\delta}}_m$ needed to calculate
$\widehat{\mathbf{S}}$ is the 2SLS estimator.

\subsection*{Three-Stage Least Squares (3SLS)}

When the set of instruments is the same across equations, the FIVE formula can
be simplified somewhat. The simplified formula is the \textbf{3SLS estimator},
denoted $\hat{\boldsymbol{\delta}}_{\mathrm{3SLS}}$.\footnote{The estimator is due to
Zellner and Theil (1962).}
To this end, let
\begin{equation}
\underset{(M \times 1)}{\boldsymbol{\varepsilon}_i}
\equiv \begin{bmatrix} \varepsilon_{i1} \\ \vdots \\ \varepsilon_{iM} \end{bmatrix}.
\label{eq:4.5.5}
\end{equation}
Then the $M \times M$ matrix of cross moments of $\boldsymbol{\varepsilon}_i$, denoted
$\boldsymbol{\Sigma}$, can be written as
\begin{equation}
\underset{(M \times M)}{\boldsymbol{\Sigma}}
\equiv \begin{bmatrix}
\sigma_{11} & \cdots & \sigma_{1M} \\
\vdots & & \vdots \\
\sigma_{M1} & \cdots & \sigma_{MM}
\end{bmatrix}
= \E(\boldsymbol{\varepsilon}_i \boldsymbol{\varepsilon}_i').
\label{eq:4.5.6}
\end{equation}

To estimate $\boldsymbol{\Sigma}$ consistently, we need an initial consistent estimator of
$\boldsymbol{\delta}_m$ for the purpose of calculating the residual $\hat{\varepsilon}_{im}$.
The term 3SLS comes from the fact that the 2SLS estimator of $\boldsymbol{\delta}_m$
is used as the initial estimator. Given the residuals thus calculated, a natural
estimator of $\boldsymbol{\Sigma}$ is
\begin{equation}
\widehat{\boldsymbol{\Sigma}}
\equiv \begin{bmatrix}
\hat{\sigma}_{11} & \cdots & \hat{\sigma}_{1M} \\
\vdots & & \vdots \\
\hat{\sigma}_{M1} & \cdots & \hat{\sigma}_{MM}
\end{bmatrix}
= \frac{1}{n} \sum_{i=1}^{n} \hat{\boldsymbol{\varepsilon}}_i\,
\hat{\boldsymbol{\varepsilon}}_i',
\label{eq:4.5.7}
\end{equation}
which is a matrix whose $(m,h)$ element is the estimated cross moment given
in~\eqref{eq:4.5.4}.

If $\mathbf{x}_i$ $(= \mathbf{x}_{i1} = \mathbf{x}_{i2} = \cdots = \mathbf{x}_{iM})$
is the common set of instruments with dimension $K$, then $\mathbf{g}_i$
in~\eqref{eq:4.1.4}, $\mathbf{S}$ in~\eqref{eq:4.5.2}, and $\widehat{\mathbf{S}}$
in~\eqref{eq:4.5.3} can be written compactly using the Kronecker
product\footnote{See formulas (A.11) and (A.15) of the Appendix if you are not
familiar with the Kronecker product notation.} as
\begin{align}
\underset{(MK \times 1)}{\mathbf{g}_i} &= \boldsymbol{\varepsilon}_i \otimes \mathbf{x}_i,
\label{eq:4.5.8} \\
\underset{(MK \times MK)}{\mathbf{S}} &= \underset{(M \times M)}{\boldsymbol{\Sigma}}
\otimes \underset{(K \times K)}{\E(\mathbf{x}_i \mathbf{x}_i')},
\label{eq:4.5.9}
\end{align}
so $\mathbf{S}^{-1} = \boldsymbol{\Sigma}^{-1} \otimes
[\E(\mathbf{x}_i \mathbf{x}_i')]^{-1}$,
\begin{equation}
\widehat{\mathbf{S}} = \widehat{\boldsymbol{\Sigma}} \otimes
\Bigl(\frac{1}{n} \sum_{i=1}^{n} \mathbf{x}_i \mathbf{x}_i'\Bigr),
\label{eq:4.5.10}
\end{equation}
so $\widehat{\mathbf{S}}^{-1} = \widehat{\boldsymbol{\Sigma}}^{-1} \otimes
\bigl(\frac{1}{n} \sum_{i=1}^{n} \mathbf{x}_i \mathbf{x}_i'\bigr)^{-1}$.
The Kronecker product decomposition~\eqref{eq:4.5.9} of $\mathbf{S}$ and
the nonsingularity of $\mathbf{S}$ (by Assumption~4.5) imply that both
$\boldsymbol{\Sigma}$ and $\E(\mathbf{x}_i \mathbf{x}_i')$ are nonsingular.

To achieve further rewriting of the 3SLS formulas, we need to go back to~\eqref{eq:4.2.6},
which, with $\hat{\mathbf{W}} = \widehat{\mathbf{S}}^{-1}$, writes out in full the
efficient multiple-equation GMM estimator. The formula~\eqref{eq:4.5.10} implies
that the $\hat{\mathbf{W}}$ in~\eqref{eq:4.2.6} is such that
\begin{equation}
\hat{\mathbf{W}}_{mh}
\;(\equiv (m,h) \text{ block of } \hat{\mathbf{W}})
= \hat{\sigma}^{mh} \cdot \Bigl(\frac{1}{n} \sum_{i=1}^{n}
\mathbf{x}_i \mathbf{x}_i'\Bigr)^{-1},
\label{eq:4.5.11}
\end{equation}
where $\hat{\sigma}^{mh}$ is the $(m,h)$ element of
$\widehat{\boldsymbol{\Sigma}}^{-1}$.
Substitute this into~\eqref{eq:4.2.6} to obtain
\begin{equation}
\hat{\boldsymbol{\delta}}_{\mathrm{3SLS}}
= \begin{bmatrix}
\hat{\sigma}^{11} \widehat{\mathbf{A}}_{11} & \cdots &
\hat{\sigma}^{1M} \widehat{\mathbf{A}}_{1M} \\
\vdots & & \vdots \\
\hat{\sigma}^{M1} \widehat{\mathbf{A}}_{M1} & \cdots &
\hat{\sigma}^{MM} \widehat{\mathbf{A}}_{MM}
\end{bmatrix}^{-1}
\begin{bmatrix}
\hat{\sigma}^{11} \hat{\mathbf{c}}_{11} + \cdots + \hat{\sigma}^{1M} \hat{\mathbf{c}}_{1M} \\
\vdots \\
\hat{\sigma}^{M1} \hat{\mathbf{c}}_{M1} + \cdots + \hat{\sigma}^{MM} \hat{\mathbf{c}}_{MM}
\end{bmatrix},
\label{eq:4.5.12}
\end{equation}
where
\begin{align}
\widehat{\mathbf{A}}_{mh}
&\equiv \Bigl(\frac{1}{n} \sum_{i=1}^{n} \mathbf{z}_{im} \mathbf{x}_i'\Bigr)
\Bigl(\frac{1}{n} \sum_{i=1}^{n} \mathbf{x}_i \mathbf{x}_i'\Bigr)^{-1}
\Bigl(\frac{1}{n} \sum_{i=1}^{n} \mathbf{x}_i \mathbf{z}_{ih}'\Bigr),
\label{eq:4.5.13} \\
\hat{\mathbf{c}}_{mh}
&\equiv \Bigl(\frac{1}{n} \sum_{i=1}^{n} \mathbf{z}_{im} \mathbf{x}_i'\Bigr)
\Bigl(\frac{1}{n} \sum_{i=1}^{n} \mathbf{x}_i \mathbf{x}_i'\Bigr)^{-1}
\Bigl(\frac{1}{n} \sum_{i=1}^{n} \mathbf{x}_i \cdot y_{ih}\Bigr).
\label{eq:4.5.14}
\end{align}

Similarly, substituting into~\eqref{eq:4.3.3} the expression for
$\boldsymbol{\Sigma}_{\mathbf{xz}}$ in~\eqref{eq:4.1.9} and the expression
for $\mathbf{S}$ in~\eqref{eq:4.5.9} and using formula (A.6) of the Appendix,
we obtain the expressions for $\avar(\hat{\boldsymbol{\delta}}_{\mathrm{3SLS}})$:
\begin{equation}
\avar(\hat{\boldsymbol{\delta}}_{\mathrm{3SLS}})
= \begin{bmatrix}
\sigma^{11} \mathbf{A}_{11} & \cdots & \sigma^{1M} \mathbf{A}_{1M} \\
\vdots & & \vdots \\
\sigma^{M1} \mathbf{A}_{M1} & \cdots & \sigma^{MM} \mathbf{A}_{MM}
\end{bmatrix}^{-1}
\label{eq:4.5.15}
\end{equation}
where
\begin{equation}
\mathbf{A}_{mh} \equiv \E(\mathbf{z}_{im} \mathbf{x}_i')
\bigl[\E(\mathbf{x}_i \mathbf{x}_i')\bigr]^{-1}
\E(\mathbf{x}_i \mathbf{z}_{ih}'),
\label{eq:4.5.16}
\end{equation}
and $\sigma^{mh}$ is the $(m,h)$ element of $\boldsymbol{\Sigma}^{-1}$.
This is consistently estimated by
\begin{equation}
\widehat{\avar}(\hat{\boldsymbol{\delta}}_{\mathrm{3SLS}})
= \begin{bmatrix}
\hat{\sigma}^{11} \widehat{\mathbf{A}}_{11} & \cdots &
\hat{\sigma}^{1M} \widehat{\mathbf{A}}_{1M} \\
\vdots & & \vdots \\
\hat{\sigma}^{M1} \widehat{\mathbf{A}}_{M1} & \cdots &
\hat{\sigma}^{MM} \widehat{\mathbf{A}}_{MM}
\end{bmatrix}^{-1}.
\label{eq:4.5.17}
\end{equation}
(This sample version can also be obtained directly by substituting~\eqref{eq:4.2.2}
and~\eqref{eq:4.5.10} into~\eqref{eq:4.3.4}.)

Most textbooks give expressions even prettier than these by using data matrices
(such as the $n \times K$ matrix $\mathbf{X}$). Deriving such expressions is an analytical
exercise to this chapter. The preceding discussion can be summarized as

\begin{proposition}[large-sample properties of 3SLS]
\label{prop:4.5}
\textit{Suppose Assumptions~4.1--4.5 and~4.7 hold, and suppose
$\mathbf{x}_{im} = \mathbf{x}_i$ (common set of instruments).
Suppose, furthermore, that\/ $\E(\mathbf{z}_{im} \mathbf{z}_{ih}')$ exists and is
finite for all $m, h$ $(= 1, 2, \ldots, M)$. Let\/ $\widehat{\boldsymbol{\Sigma}}$
be the $M \times M$ matrix of estimated error cross moments calculated
by~\eqref{eq:4.5.7} using the 2SLS residuals. Then}
\begin{enumerate}[label=(\alph*)]
\item $\hat{\boldsymbol{\delta}}_{\mathrm{3SLS}}$ \textit{given by~\eqref{eq:4.5.12}
is consistent, asymptotically normal, and efficient, with\/
$\avar(\hat{\boldsymbol{\delta}}_{\mathrm{3SLS}})$ given by~\eqref{eq:4.5.15}.}

\item \textit{The estimated asymptotic variance~\eqref{eq:4.5.17} is consistent for\/
$\avar(\hat{\boldsymbol{\delta}}_{\mathrm{3SLS}})$.}

\item \textit{(Sargan's statistic)}
\[
J(\hat{\boldsymbol{\delta}}_{\mathrm{3SLS}},\, \widehat{\mathbf{S}}^{-1})
\equiv n \cdot \mathbf{g}_n(\hat{\boldsymbol{\delta}}_{\mathrm{3SLS}})'
\widehat{\mathbf{S}}^{-1}\, \mathbf{g}_n(\hat{\boldsymbol{\delta}}_{\mathrm{3SLS}})
\dto \chi^2\!\Bigl(MK - \sum_m L_m\Bigr),
\]
\textit{where\/ $\widehat{\mathbf{S}} = \widehat{\boldsymbol{\Sigma}} \otimes
\bigl(\frac{1}{n} \sum_i \mathbf{x}_i \mathbf{x}_i'\bigr)$,
$K$ is the number of common instruments, and\/ $\mathbf{g}_n(\cdot)$ is given
in~\eqref{eq:4.2.1}.}
\end{enumerate}
\end{proposition}

\subsection*{Seemingly Unrelated Regressions (SUR)}

The 3SLS formula can further be simplified if
\begin{equation}
\mathbf{x}_i = \text{union of } (\mathbf{z}_{i1}, \ldots, \mathbf{z}_{iM}).
\label{eq:4.5.18}
\end{equation}
This is equivalent to the condition that
\begin{equation}
\E(\mathbf{z}_{im} \cdot \varepsilon_{ih}) = \mathbf{0}
\quad (m, h = 1, 2, \ldots, M).
\tag{4.5.18$'$}
\label{eq:4.5.18p}
\end{equation}
That is, the predetermined regressors satisfy ``cross'' orthogonalities: not only are
they predetermined in each equation (i.e., $\E(\mathbf{z}_{im} \cdot \varepsilon_{im}) = \mathbf{0}$),
but also they are predetermined in the other equations (i.e.,
$\E(\mathbf{z}_{im} \cdot \varepsilon_{ih}) = \mathbf{0}$ for $m \neq h$).
The simplified formula is called the \textbf{SUR estimator}, to be denoted
$\hat{\boldsymbol{\delta}}_{\mathrm{SUR}}$.\footnote{The estimator is due to Zellner (1962).}

\begin{quote}
\textbf{Example 4.3 (Example 4.1 with different specification for instruments):}\quad
If~\eqref{eq:4.5.18} holds for the two-equation system of Example~4.1,
$\mathbf{x}_i$ is the union of $\mathbf{z}_{i1}$ and $\mathbf{z}_{i2}$, so
\begin{equation}
\mathbf{x}_i = \begin{bmatrix} 1 \\ S_i \\ IQ_i \\ EXPR_i \end{bmatrix}.
\label{eq:4.5.19}
\end{equation}
The orthogonality conditions that $\E(\mathbf{x}_i \cdot \varepsilon_{i1}) = \mathbf{0}$
and $\E(\mathbf{x}_i \cdot \varepsilon_{i2}) = \mathbf{0}$ become
\begin{equation}
\E \begin{bmatrix}
\varepsilon_{i1} \\ S_i\, \varepsilon_{i1} \\
IQ_i\, \varepsilon_{i1} \\ EXPR_i\, \varepsilon_{i1}
\end{bmatrix} = \mathbf{0}
\quad \text{and} \quad
\E \begin{bmatrix}
\varepsilon_{i2} \\ S_i\, \varepsilon_{i2} \\
IQ_i\, \varepsilon_{i2} \\ EXPR_i\, \varepsilon_{i2}
\end{bmatrix} = \mathbf{0},
\label{eq:4.5.20}
\end{equation}
which should be contrasted with the statement that the regressors are predetermined
in each equation:
\begin{equation}
\E \begin{bmatrix}
\varepsilon_{i1} \\ S_i\, \varepsilon_{i1} \\
IQ_i\, \varepsilon_{i1} \\ EXPR_i\, \varepsilon_{i1}
\end{bmatrix} = \mathbf{0}
\quad \text{and} \quad
\E \begin{bmatrix}
\varepsilon_{i2} \\ S_i\, \varepsilon_{i2} \\
IQ_i\, \varepsilon_{i2}
\end{bmatrix} = \mathbf{0}.
\label{eq:4.5.21}
\end{equation}
The difference between~\eqref{eq:4.5.20} and~\eqref{eq:4.5.21} is that the $KWW$
equation is overidentified with $EXPR$ as the additional instrument.
\end{quote}

Because SUR is a special case of 3SLS, formulas~\eqref{eq:4.5.12},
\eqref{eq:4.5.15}, and~\eqref{eq:4.5.17} for $\hat{\boldsymbol{\delta}}_{\mathrm{3SLS}}$
apply to $\hat{\boldsymbol{\delta}}_{\mathrm{SUR}}$.
The implication of the SUR assumption~\eqref{eq:4.5.18} is that, in
the above expressions for $\widehat{\mathbf{A}}_{mh}$, $\hat{\mathbf{c}}_{mh}$,
and $\mathbf{A}_{mh}$, $\mathbf{x}_i$ ``disappears'':
\begin{align}
\widehat{\mathbf{A}}_{mh} &= \frac{1}{n} \sum_{i=1}^{n}
\mathbf{z}_{im} \mathbf{z}_{ih}',
\tag{4.5.13$'$} \label{eq:4.5.13p} \\
\hat{\mathbf{c}}_{mh} &= \frac{1}{n} \sum_{i=1}^{n}
\mathbf{z}_{im} \cdot y_{ih},
\tag{4.5.14$'$} \label{eq:4.5.14p} \\
\mathbf{A}_{mh} &= \E(\mathbf{z}_{im} \mathbf{z}_{ih}').
\tag{4.5.16$'$} \label{eq:4.5.16p}
\end{align}

Here we give a proof of~\eqref{eq:4.5.16p}. Without loss of generality, suppose
$\mathbf{z}_{im}$ is the first $L_m$ elements of the $K$ elements of $\mathbf{x}_i$.
Let $\mathbf{D}$ $(K \times L_m)$ be the first $L_m$ columns of $\mathbf{I}_K$.
Then we have
\[
\mathbf{z}_{im} = \mathbf{D}' \mathbf{x}_i, \quad
\E(\mathbf{z}_{im} \mathbf{x}_i') = \mathbf{D}'\, \E(\mathbf{x}_i \mathbf{x}_i'), \quad
\mathbf{D}\, \E(\mathbf{x}_i \mathbf{z}_{ih}') = \E(\mathbf{z}_{im} \mathbf{z}_{ih}').
\]
Substituting these into~\eqref{eq:4.5.16} yields~\eqref{eq:4.5.16p}.
The proof of~\eqref{eq:4.5.13p} and~\eqref{eq:4.5.14p}, which are about the sample
moments, proceeds similarly, with $\mathbf{z}_{im} = \mathbf{D}' \mathbf{x}_i$.
(You will be asked to provide an alternative proof by the use of the projection matrix
in the problem set.)

In 3SLS, the initial consistent estimator for calculating
$\widehat{\boldsymbol{\Sigma}}$ was 2SLS. But, 2SLS when the regressors are a subset
of the instrument set is OLS (as you showed in Review Question~7 of Section~3.8).
So, for SUR, the initial estimator is the OLS estimator. The preceding discussion
can be summarized as

\begin{proposition}[large-sample properties of SUR]
\label{prop:4.6}
\textit{Suppose Assumptions~4.1--4.5 and~4.7 hold with
$\mathbf{x}_i = \text{union of}\; (\mathbf{z}_{i1}, \ldots, \mathbf{z}_{iM})$.
Let\/ $\widehat{\boldsymbol{\Sigma}}$ be the $M \times M$ matrix of estimated error
cross moments calculated by~\eqref{eq:4.5.7} using the OLS residuals. Then}
\begin{enumerate}[label=(\alph*)]
\item $\hat{\boldsymbol{\delta}}_{\mathrm{SUR}}$ \textit{given by~\eqref{eq:4.5.12}
with\/ $\widehat{\mathbf{A}}_{mh}$ and\/ $\hat{\mathbf{c}}_{mh}$ given
by~\eqref{eq:4.5.13p} and~\eqref{eq:4.5.14p} for
$m, h = 1, \ldots, M$ is consistent, asymptotically normal, and efficient, with\/
$\avar(\hat{\boldsymbol{\delta}}_{\mathrm{SUR}})$ given by~\eqref{eq:4.5.15}
where\/ $\mathbf{A}_{mh}$ is given by~\eqref{eq:4.5.16p}.}

\item \textit{The estimated asymptotic variance~\eqref{eq:4.5.17} where\/
$\widehat{\mathbf{A}}_{mh}$ is given by~\eqref{eq:4.5.13p} is consistent for\/
$\avar(\hat{\boldsymbol{\delta}}_{\mathrm{SUR}})$.}

\item \textit{(Sargan's statistic)}
\[
J(\hat{\boldsymbol{\delta}}_{\mathrm{SUR}},\, \widehat{\mathbf{S}}^{-1})
\equiv n \cdot \mathbf{g}_n(\hat{\boldsymbol{\delta}}_{\mathrm{SUR}})'
\widehat{\mathbf{S}}^{-1}\, \mathbf{g}_n(\hat{\boldsymbol{\delta}}_{\mathrm{SUR}})
\dto \chi^2\!\Bigl(MK - \sum_m L_m\Bigr),
\]
\textit{where\/ $\widehat{\mathbf{S}} = \widehat{\boldsymbol{\Sigma}} \otimes
\bigl(\frac{1}{n} \sum_i \mathbf{x}_i \mathbf{x}_i'\bigr)$,
$K$ is the number of common instruments, and\/ $\mathbf{g}_n(\cdot)$ is given
in~\eqref{eq:4.2.1}.}
\end{enumerate}
\end{proposition}

(Unlike in Proposition~4.5, the condition that $\E(\mathbf{z}_{im} \mathbf{z}_{ih}')$
be finite is not needed because, thanks to~\eqref{eq:4.5.18}, it is implied by the
condition that $\E(\mathbf{x}_i \mathbf{x}_i')$ be nonsingular, which in turn is implied
by the nonsingularity of $\mathbf{S}$ [from Assumption~4.5] and the fact that
$\mathbf{S}$ can be written as $\boldsymbol{\Sigma} \otimes \E(\mathbf{x}_i \mathbf{x}_i')$.)
Sargan's statistic for SUR is rarely reported in practice.

\subsection*{SUR versus OLS}

Since the regressors are predetermined, the system can also be estimated by the
equation-by-equation OLS. Then why SUR over OLS? As just seen, under conditional
homoskedasticity, the efficient multiple-equation GMM is FIVE, which in
turn is numerically equal to SUR under~\eqref{eq:4.5.18}. As shown in Chapter~3,
under conditional homoskedasticity, the efficient single-equation GMM is 2SLS,
which in turn is numerically equal to OLS because the regressors are predetermined
under~\eqref{eq:4.5.18}. This interrelationship is shown in Figure~4.1.
Therefore, the relation between SUR and equation-by-equation OLS is strictly
analogous to the relation between the multiple-equation GMM and the
equation-by-equation GMM. As discussed in Section~4.4, there are two cases to consider.

\begin{enumerate}[label=(\alph*)]
\item Each equation is just identified. Because the common instrument set is the
union of all the regressors, this is possible only if the regressors are the same
for all equations, i.e., if $\mathbf{z}_{im} = \mathbf{x}_i$ for all $m$.
The SUR for this case is called the \textbf{multivariate regression}.
We observed in Section~4.4 that both the efficient multiple-equation GMM and
the efficient single-equation GMM are numerically the same as the IV estimator
for the just identified system. Because the regressors are predetermined, we
conclude that the GMM estimator of the multivariate regression model is simply
equation-by-equation OLS.\footnote{It will be shown in Chapter~8 that the
estimator is also the maximum likelihood estimator.}
This can be verified directly by substituting~\eqref{eq:4.5.13p}
and~\eqref{eq:4.5.14p} on page~280 (with $\mathbf{z}_{im} = \mathbf{x}_i$)
into the expression for the point estimate~\eqref{eq:4.5.12}.
Also substituting (again with $\mathbf{z}_{im} = \mathbf{x}_i$)
into~\eqref{eq:4.5.16p} and~\eqref{eq:4.5.13p}
into~\eqref{eq:4.5.15} and~\eqref{eq:4.5.17}, we obtain, for multivariate regression,
\begin{align}
\avar(\hat{\boldsymbol{\delta}})
&= \boldsymbol{\Sigma} \otimes \bigl[\E(\mathbf{x}_i \mathbf{x}_i')\bigr]^{-1},
\label{eq:4.5.22} \\
\widehat{\avar}(\hat{\boldsymbol{\delta}})
&= \widehat{\boldsymbol{\Sigma}} \otimes
\Bigl(\frac{1}{n} \sum_{i=1}^{n} \mathbf{x}_i \mathbf{x}_i'\Bigr)^{-1}.
\label{eq:4.5.23}
\end{align}

\item At least one equation is overidentified. Then SUR is more efficient than
equation-by-equation OLS, unless equations are ``unrelated'' to each other in
the sense of~\eqref{eq:4.4.3}. In the present case of conditional homoskedasticity
and the common set of instruments,~\eqref{eq:4.4.3} becomes
\[
\sigma_{mh}\, \E(\mathbf{x}_i \mathbf{x}_i') = \mathbf{0}
\quad \text{for all } m \neq h.
\]
Since $\E(\mathbf{x}_i \mathbf{x}_i')$ cannot be a zero matrix (it is assumed to be
nonsingular), equations are ``unrelated'' to each other if and only if
$\sigma_{mh} = 0$. Therefore, SUR is more efficient than OLS if $\sigma_{mh} \neq 0$
for some pair $(m, h)$. If $\sigma_{mh} = 0$ for all $(m, h)$, the two estimators are
asymptotically equivalent (i.e., $\sqrt{n}$ times the difference vanishes).
\end{enumerate}

Another way to see the efficiency of SUR is to view the SUR model as
a multivariate regression model with \textit{a priori} exclusion restrictions. As an
example, expand the two-equation system of Example~4.3 as
\begin{align*}
LW_i &= \phi_1 + \beta_1 S_i + \gamma_1 IQ_i + \pi_1 EXPR_i + \varepsilon_{i1}, \\
KWW_i &= \phi_2 + \beta_2 S_i + \gamma_2 IQ_i + \pi_2 EXPR_i + \varepsilon_{i2}.
\end{align*}
The common instrument set is $(1, S_i, IQ_i, EXPR_i)$. Then this two-equation system
is a multivariate regression model with the same regressors in both equations.
But if $\pi_2$ is \textit{a priori} restricted to be $0$ and thus $EXPR_i$ is excluded from
the second equation, then the model becomes the SUR model. The SUR estimator
is more efficient than the multivariate regression because it exploits exclusion
restrictions.

The relationship among the estimators considered in this section is summarized
in Table~4.2.

% Figure 4.1 placeholder
\begin{figure}[htbp]
\centering
% TODO: Create figure showing relationship between OLS, 2SLS, SUR, 3SLS, FIVE, and GMM
\includegraphics[width=0.9\textwidth,trim=50 350 50 350,clip]{figures/fig_4_1.png}
\caption{OLS and GMM}
\label{fig:4.1}
\end{figure}

% Table 4.2
\begin{table}[htbp]
\centering
\caption{Relationship between Multiple-Equation Estimators}
\label{tab:4.2}
\small
\medskip
\begin{tabular}{@{}p{2cm}p{1.6cm}p{1.6cm}p{2cm}p{1.8cm}p{2cm}@{}}
\toprule
& Multiple-eq.\ GMM & FIVE & 3SLS & SUR & Multivariate regression \\
\midrule
The model
& Ass.\ 4.1--4.6
& Ass.\ 4.1--4.5, 4.7, $\E(\mathbf{z}_{im}\mathbf{z}_{ih}')$ finite
& Ass.\ 4.1--4.5, 4.7, $\E(\mathbf{z}_{im}\mathbf{z}_{ih}')$ finite,
$\mathbf{x}_{im} = \mathbf{x}_i$ for all $m$
& Ass.\ 4.1--4.5, 4.7, $\mathbf{x}_{im} = \mathbf{x}_i$ for all $m$,
$\mathbf{x}_i = \text{union of } (\mathbf{z}_{i1}, \ldots, \mathbf{z}_{iM})$
& Ass.\ 4.1--4.5, 4.7, $\mathbf{x}_{im} = \mathbf{x}_i$ for all $m$,
$\mathbf{z}_{im} = \mathbf{x}_i$ for all $m$ \\
\midrule
$\mathbf{S}$ & (4.1.11) & (4.5.2) & \multicolumn{3}{l}{$\boldsymbol{\Sigma} \otimes \E(\mathbf{x}_i \mathbf{x}_i')$} \\[3pt]
$\widehat{\mathbf{S}}$ & (4.3.2) & (4.5.3) & \multicolumn{3}{l}{$\widehat{\boldsymbol{\Sigma}} \otimes (n^{-1} \sum_i \mathbf{x}_i \mathbf{x}_i')$} \\[3pt]
$\hat{\boldsymbol{\delta}}(\widehat{\mathbf{S}}^{-1})$
& no simplif.
& no simplif.
& (4.5.12) with (4.5.13), (4.5.14)
& (4.5.12) with (4.5.13$'$), (4.5.14$'$)
& eq-by-eq OLS \\
\bottomrule
\end{tabular}
\end{table}
